%! Zeros of Polynomials: the Rational Root Theorem — Unit 3, Lesson 3.4 — Student Packet Cover Sheet
\documentclass[10pt]{article}
\usepackage{algebra2-article}
\usepackage{algebra2-boxes}
\usepackage{ltablex}
\keepXColumns

\begin{document}

% Full-bleed forest banner
\begin{tikzpicture}[remember picture, overlay]
  \fill[forest] (current page.north west)
    rectangle ([yshift=-0.9in]current page.north east);
\end{tikzpicture}
\vspace{-0.8in}

\begin{center}
  {\color{white}\textbf{\LARGE Algebra 2}}\\[4pt]
  {\color{forestmid}\large\textbf{Unit 3 \quad Polynomial Functions}}\\[6pt]
  {\color{white}\textbf{\large Lesson 3.4 \quad Zeros of Polynomials: the Rational Root Theorem}}
\end{center}

\vspace{0.22in}
\namedateperiod
\vspace{0.18in}

\begin{learningtargetbox}
\small
\begin{enumerate}[leftmargin=1.5em, itemsep=5pt]
  \item \ldots{} use the \textbf{Rational Root Theorem} to list every possible rational zero of a polynomial with integer coefficients: $\pm\frac{p}{q}$, with $p$ a divisor of the \textbf{constant term} and $q$ a divisor of the \textbf{leading coefficient}.
  \item \ldots{} reduce that list and keep the $\pm$ --- no repeats, no missing negatives.
  \item \ldots{} \textbf{test} candidates efficiently using the Remainder Theorem, synthetic division, and the $x$-intercepts of a graph.
  \item \ldots{} \textbf{divide} out each zero I find and factor the depressed polynomial \textbf{completely}.
  \item \ldots{} solve a polynomial equation of degree $3$ or higher and list all of its real solutions.
  \item \ldots{} \textbf{verify} my solutions by substitution and against a graph, and explain why a candidate list can come up empty.
\end{enumerate}
\end{learningtargetbox}

\vspace{0.14in}

\begin{tocbox}
\small
\renewcommand{\arraystretch}{1.5}
\begin{tabularx}{\linewidth}{@{} c l X r @{}}
\rowcolor{forestbg}
\textbf{\#} & \textbf{Component} & \textbf{Description} & \textbf{Score} \\
\midrule
1 & Warm-Up        & Divisors, building fractions, and finishing yesterday's polynomial      & \blank{1.2cm} \\
2 & Guided Notes   & The Rational Root Theorem and the list--test--divide--finish workflow   & \blank{1.2cm} \\
3 & Group Activity & \emph{Suspect List} --- hunting rational zeros, by tier                 & \blank{1.2cm} \\
4 & Exit Ticket    & Quick check: one full solve, one error analysis, one SOL-style item     & \blank{1.2cm} \\
5 & Homework       & Candidate lists, cubics and a quartic, and a planter-box model          & \blank{1.2cm} \\
\midrule
  & \multicolumn{2}{r}{\textbf{Total}} & \blank{1.2cm} \\
\end{tabularx}
\end{tocbox}

\vspace{0.10in}

\begin{remindbox}
\small Lesson 3.3 gave you everything you needed to crack a polynomial --- \emph{once somebody told
you which number to test}. Today you build that list yourself. If $\frac{p}{q}$ is a rational zero of
a polynomial with integer coefficients, then $p$ divides the \textbf{constant term} and $q$ divides
the \textbf{leading coefficient}: constant on \textbf{top}, leading coefficient on the
\textbf{bottom}, every value with a $\pm$. That turns infinitely many numbers into a short list you
can actually test. Two cautions come with it. The list holds \emph{possible} zeros, not guaranteed
ones --- a polynomial can have none of them. And the theorem finds \textbf{rational} zeros only;
irrational zeros like $\pm\sqrt{3}$ show up in the depressed polynomial, never on the list.
\end{remindbox}

\end{document}
